\documentclass[11pt]{article}

\usepackage[margin=1in]{geometry}
\usepackage{amsmath, amssymb}
\usepackage[T1]{fontenc}
\usepackage{lmodern}
\usepackage{xcolor}
\usepackage{hyperref}
\hypersetup{
  colorlinks=true,
  linkcolor=blue,
  urlcolor=blue,
  citecolor=blue
}
\usepackage{enumitem}
\setlist{nosep}
\usepackage{listings}
\lstset{
  basicstyle=\ttfamily\small,
  columns=fullflexible,
  frame=single,
  breaklines=true
}

\title{MC Dropout \& Bayesian Uncertainty in the Muse 2 EEG Finger/Action Pipeline}
\author{Repository-local implementation notes}
\date{\today}

\begin{document}
\maketitle

\section{Goal and context (why uncertainty exists in this pipeline)}

This repository builds an end-to-end EEG decoding pipeline for predicting:
(i) an \textbf{action} state (REST / OPEN / CLOSE), and
(ii) a \textbf{finger} identity (NONE / THUMB / INDEX / MIDDLE / RING / PINKY).

EEG is noisy, non-stationary, and susceptible to artifacts and distribution shift (e.g., electrode movement, muscle noise, attention changes).
Because downstream use can be safety-relevant (e.g., actuation), the pipeline needs not just a class prediction but also an estimate of \textbf{how reliable} that prediction is.

This repo implements \textbf{Monte Carlo (MC) dropout} as a practical Bayesian approximation to estimate \textbf{model uncertainty} at inference time without changing the trained architecture.

\section{Where MC dropout and uncertainty live in this repository (code map)}

The following files are the main implementation points:

\begin{itemize}
  \item \textbf{Dropout layers (the stochasticity we sample).}
  The deployed model class is \texttt{CNNLSTMFingerActionNet} in \texttt{models/cnn\_lstm\_finger\_action\_net.py}.
  Dropout is applied in the convolutional stack (\texttt{nn.Dropout(p=dropout)} at \texttt{models/cnn\_lstm\_finger\_action\_net.py:39}) and again before the heads (\texttt{self.head\_dropout} at \texttt{models/cnn\_lstm\_finger\_action\_net.py:50}, used at \texttt{models/cnn\_lstm\_finger\_action\_net.py:62}).

  \item \textbf{MC-dropout forward pass (core uncertainty computation).}
  \texttt{CNNLSTMFingerActionNet.mc\_forward} at \texttt{models/cnn\_lstm\_finger\_action\_net.py:89} enables dropout via \texttt{self.train()} (\texttt{models/cnn\_lstm\_finger\_action\_net.py:93}), runs \texttt{passes} stochastic forwards, and returns per-head mean/std plus additional uncertainty metrics (entropy and mutual information).

  \item \textbf{Online inference wrapper + safety gating.}
  \texttt{utils/inference.py} provides \texttt{InferenceEngine} (\texttt{utils/inference.py:54}).
  It computes scalar uncertainties from MC dropout (\texttt{utils/inference.py:185}) and uses them for an adaptive actuation gate (\texttt{utils/inference.py:221}).

  \item \textbf{Offline uncertainty visualization for paper-style figures.}
  \texttt{3c\_live\_paper\_figures.py} explicitly runs MC dropout with \texttt{model.mc\_forward} (\texttt{3c\_live\_paper\_figures.py:232}) and computes per-sample scalar uncertainty as the mean of per-class MC std (\texttt{3c\_live\_paper\_figures.py:242} and \texttt{3c\_live\_paper\_figures.py:246}).

  \item \textbf{Legacy MC-dropout helper (not wired into the main pipeline).}
  \texttt{utils/mc\_dropout.py:6} contains \texttt{mc\_dropout\_predict}, which forces \texttt{model.train()} to enable dropout and returns mean/std.
  In this repo state, the \emph{main} path uses \texttt{CNNLSTMFingerActionNet.mc\_forward} and/or \texttt{utils/inference.\_mc\_predict}.

  \item \textbf{Configuration/documentation knobs.}
  \texttt{docs/meta/Project Structure.ini} records intent (\texttt{bayesian\_mode = mc\_dropout}, \texttt{uncertainty\_measure = mc\_dropout\_std}, \texttt{mc\_samples = 30}; see \texttt{docs/meta/Project Structure.ini:41} and \texttt{docs/meta/Project Structure.ini:49}).
  The GUI lists an \texttt{MC\_DROPOUT\_PASSES} setting in \texttt{eeglab\_wrapper\_ui.py:395} and shows it in a dropdown (\texttt{eeglab\_wrapper\_ui.py:2996}).
\end{itemize}

\section{MC dropout as a Bayesian approximation (the underlying idea)}

\subsection{Deterministic prediction}
A neural classifier with weights $W$ produces logits $z=f_W(x)$ and probabilities via softmax:
\[
p(y=c\mid x, W)=\mathrm{softmax}(z)_c=\frac{e^{z_c}}{\sum_{j=1}^K e^{z_j}}.
\]
In standard inference you call \texttt{model.eval()}, which disables dropout and yields a single probability vector for each input.

\subsection{Dropout introduces stochastic subnetworks}
Dropout multiplies activations by a random Bernoulli mask.
Conceptually, for a hidden unit $h_i$:
\[
m_i \sim \mathrm{Bernoulli}(1-p),\qquad \tilde{h}_i = m_i\cdot h_i.
\]
Each dropout mask corresponds to a different ``thinned'' subnetwork, so repeating inference with dropout enabled yields a \emph{distribution} of predictions.

\subsection{Bayesian predictive distribution and MC approximation}
The Bayesian posterior predictive distribution marginalizes weights:
\[
p(y\mid x,\mathcal{D})=\int p(y\mid x,W)\,p(W\mid \mathcal{D})\,dW.
\]
MC dropout approximates this integral by sampling $T$ dropout masks (i.e., $W^{(t)}$ induced by masks) and averaging:
\[
p(y\mid x,\mathcal{D})\approx \frac{1}{T}\sum_{t=1}^T p(y\mid x, W^{(t)}).
\]
This is the core operational definition used in this repository.

\section{How uncertainty is computed in this repo (exact computations)}

Let $T$ be the number of MC passes (called \texttt{passes} or \texttt{mc\_passes} in code).
For a given input window $x$, each stochastic pass produces probabilities:
\[
p^{(t)} = p(y\mid x, W^{(t)}),\qquad t=1,\dots,T.
\]

\subsection{Predictive mean (the probability we use as the final distribution)}
\[
\bar{p} = \frac{1}{T}\sum_{t=1}^T p^{(t)}.
\]
Implementation: \texttt{finger\_mean} / \texttt{action\_mean} in
\texttt{models/cnn\_lstm\_finger\_action\_net.py:104} and
\texttt{utils/inference.py:41}.

\subsection{Per-class predictive standard deviation (disagreement across passes)}
For class $c$:
\[
\sigma_c = \sqrt{\frac{1}{T-1}\sum_{t=1}^T \left(p^{(t)}_c - \bar{p}_c\right)^2}.
\]
Implementation: \texttt{finger\_std} / \texttt{action\_std} in
\texttt{models/cnn\_lstm\_finger\_action\_net.py:106} and
\texttt{utils/inference.py:43}.

\subsection{Scalar uncertainty used downstream (mean over classwise std)}
The online engine reduces the std vector to a scalar by averaging across classes:
\[
u_{\text{action}} = \frac{1}{K_a}\sum_{c=1}^{K_a}\sigma^{\text{action}}_c,\qquad
u_{\text{finger}} = \frac{1}{K_f}\sum_{c=1}^{K_f}\sigma^{\text{finger}}_c.
\]
Implementation:
\texttt{utils/inference.py:185} and \texttt{utils/inference.py:186}.

Offline figures compute the same quantity per sample:
\texttt{3c\_live\_paper\_figures.py:242} and \texttt{3c\_live\_paper\_figures.py:246}.

\subsection{Confidence (distinct from uncertainty)}
The system also computes a conventional confidence score from the predictive mean:
\[
\mathrm{confidence}=\max_c \bar{p}_c.
\]
Implementation: \texttt{utils/inference.py:218} (action) and \texttt{utils/inference.py:219} (finger).
This is conceptually different from $u$:
\begin{itemize}
  \item Confidence is ``how large is the top probability under $\bar{p}$?''
  \item MC-dropout uncertainty is ``how much do sampled subnetworks disagree?''
\end{itemize}

\section{Additional Bayesian-style uncertainty metrics computed by the model}

The model's \texttt{mc\_forward} returns more than mean/std:

\subsection{Predictive entropy (total uncertainty)}
\[
H(\bar{p})=-\sum_{c}\bar{p}_c\log(\bar{p}_c).
\]
Implementation: \texttt{models/cnn\_lstm\_finger\_action\_net.py:109} and \texttt{:110}.

\subsection{Expected entropy (average per-sample ambiguity)}
\[
\mathbb{E}[H(p^{(t)})]\approx \frac{1}{T}\sum_{t=1}^T\left(-\sum_c p^{(t)}_c\log(p^{(t)}_c)\right).
\]
Implementation: \texttt{models/cnn\_lstm\_finger\_action\_net.py:111}--\texttt{:116}.

\subsection{Mutual information (model disagreement / epistemic proxy)}
\[
\mathrm{MI}\approx H(\bar{p})-\mathbb{E}[H(p^{(t)})].
\]
Implementation: \texttt{models/cnn\_lstm\_finger\_action\_net.py:118} and \texttt{:119}.

\paragraph{Important note.}
Even though entropy/MI are computed, the current online gating logic uses the \textbf{mean std} scalar $u$ (\texttt{utils/inference.py:185}).

\section{How uncertainty is used (online safety gate + suppression)}

The repository's explicit uncertainty-aware decision layer is implemented in \texttt{InferenceEngine.predict} (\texttt{utils/inference.py:199}).

\subsection{Adaptive confidence threshold increases with uncertainty}
Let $\tau_{\text{base}}$ be the base threshold and $\lambda$ be an uncertainty weight.
The engine computes an adaptive threshold:
\[
\tau_{\text{adaptive}} = \mathrm{clip}\left(\tau_{\text{base}} + \lambda\cdot u_{\text{action}}\right),
\]
where clipping bounds it (in code) to at most $0.99$ and at least $\tau_{\text{base}}$.

Implementation: \texttt{utils/inference.py:221}--\texttt{:228}.

\subsection{Temporal stability / debounce}
Predictions must be stable across \texttt{stability\_frames} windows before allowing actuation.
Implementation: \texttt{utils/inference.py:230}--\texttt{:234}.

\subsection{Velocity/strength suppression by uncertainty}
When action is non-REST, the engine suppresses output magnitude:
\[
\mathrm{velocity}=\mathrm{confidence}\cdot(1-u_{\text{action}}).
\]
Implementation: \texttt{utils/inference.py:236}--\texttt{:239}.

\section{Caveats and implementation details that matter}

\subsection{MC dropout requires dropout enabled at inference time}
MC dropout only works if dropout is active.
This is why code explicitly switches the module to training mode inside MC prediction:
\texttt{self.train()} in \texttt{models/cnn\_lstm\_finger\_action\_net.py:93} and
\texttt{model.train()} in \texttt{utils/inference.py:27}.

\subsection{But \texttt{train()} affects more than dropout}
\begin{itemize}
  \item \textbf{BatchNorm warning.} If a model contains BatchNorm, \texttt{train()} changes behavior (and can update running stats).
  \item \textbf{Why this repo's main model is safer.} \texttt{CNNLSTMFingerActionNet} uses GroupNorm (\texttt{models/cnn\_lstm\_finger\_action\_net.py:34}), which does not maintain running statistics, so turning on dropout via \texttt{train()} is typically acceptable for MC sampling.
  \item \textbf{Older model variant.} \texttt{models/finger\_action\_net.py} uses BatchNorm (\texttt{models/finger\_action\_net.py:33}); using MC dropout by toggling \texttt{train()} in that model would require extra care.
\end{itemize}

\subsection{Compute/latency tradeoff}
MC passes increase inference cost roughly linearly with $T$.
The online defaults are small (\texttt{InferenceConfig.mc\_passes=10} at \texttt{utils/inference.py:13}), while the paper figure script uses \texttt{MC\_SAMPLES} (see \texttt{docs/meta/Project Structure.ini:50} for the intended \texttt{mc\_samples=30}).

\section{How (and where) this is currently wired into the pipeline}

\subsection{Offline / reporting usage (wired)}
\begin{itemize}
  \item \textbf{Training (Step 2):} \texttt{2\_train\_model.py} trains \texttt{CNNLSTMFingerActionNet} (see model creation at \texttt{2\_train\_model.py:927}) and saves weights.
  \item \textbf{Deterministic evaluation (Step 3):} \texttt{3\_evaluate\_model.py} is explicitly labeled ``Deterministic evaluation --- Dropout OFF'' (\texttt{3\_evaluate\_model.py:1}).
  \item \textbf{Uncertainty figures (Step 3c):} \texttt{3c\_live\_paper\_figures.py} runs MC dropout and plots confidence vs uncertainty and reliability style plots. This is the primary ``uncertainty in evaluation'' path in the current repo.
\end{itemize}

\subsection{Online inference engine (implemented, but not the default live script)}
\begin{itemize}
  \item \texttt{InferenceEngine} is fully implemented (\texttt{utils/inference.py:54}) and includes uncertainty-aware gating (\texttt{utils/inference.py:221}).
  \item In this repo state, the only direct consumer is a profiling tool: \texttt{tools/profile\_live\_loop.py} imports and uses \texttt{InferenceEngine} (\texttt{tools/profile\_live\_loop.py:13} and \texttt{:73}).
\end{itemize}

\subsection{Step 1 and Step 7 wiring status (integration gap)}
\begin{itemize}
  \item \textbf{Step 1 is record-only.}
  \texttt{1\_stream\_and\_record.py} focuses on LSL acquisition, live plotting, and writing raw/events artifacts; it does not run the classifier and therefore does not produce MC-dropout uncertainty.
  \item \textbf{Step 7 is deterministic.}
  \texttt{7\_live\_infer\_and\_actuate.py} performs a single forward pass under \texttt{torch.inference\_mode()} and uses softmax probabilities (\texttt{7\_live\_infer\_and\_actuate.py:418}--\texttt{:421}).
  It does not call \texttt{mc\_forward} or \texttt{InferenceEngine}, and it does not compute or log uncertainty.
  \item \textbf{Schema intent vs current wiring.}
  \texttt{docs/spec/SCHEMAS.md} defines fields such as \texttt{action\_uncertainty} and \texttt{finger\_uncertainty} for feature streams (\texttt{docs/spec/SCHEMAS.md:22}), but in the current repo state, those fields are only produced by \texttt{InferenceEngine} and by the offline MC figure script, not by the live Step 7 loop.
\end{itemize}

\section{Recommended integration point (to make live uncertainty end-to-end)}

If the goal is to compute uncertainty during live inference and use it to gate actuation, the most direct integration point is:

\begin{enumerate}
  \item Replace the single deterministic forward pass in \texttt{7\_live\_infer\_and\_actuate.py} with a call to \texttt{InferenceEngine.predict} or \texttt{InferenceEngine.predict\_proba}.
  \item Feed each extracted window (shape $[T,C]$) directly to the engine; it already performs normalization and MC sampling (\texttt{utils/inference.py:110}--\texttt{:137} and \texttt{utils/inference.py:164}--\texttt{:187}).
  \item Use the engine's returned \texttt{adaptive\_threshold}, \texttt{allow\_actuation}, and uncertainty-suppressed \texttt{velocity} (\texttt{utils/inference.py:251}--\texttt{:258}) instead of the current ``min joint prob'' heuristic.
  \item Log \texttt{action\_confidence}, \texttt{action\_uncertainty}, \texttt{finger\_confidence}, \texttt{finger\_uncertainty} so downstream review/reporting matches \texttt{docs/spec/SCHEMAS.md}.
\end{enumerate}

\section{Minimal pseudocode for the intended online uncertainty-aware loop}

\begin{lstlisting}
# window_TxC: np.ndarray shape [T, C]
pred, safety, diag = engine.predict(window_TxC)

# pred contains:
#   action_id, finger_id
#   action_confidence, action_uncertainty
#   finger_confidence, finger_uncertainty
#
# safety contains:
#   adaptive_threshold, allow_actuation, velocity, stability_ok, ...

if safety["allow_actuation"]:
    send_to_robot(pred["finger_id"], pred["action_id"], safety["velocity"])
else:
    do_nothing()
\end{lstlisting}

\section{Summary}

\begin{itemize}
  \item MC dropout uncertainty is implemented in the model (\texttt{mc\_forward}) and in an online wrapper (\texttt{InferenceEngine}).
  \item The scalar uncertainty used for gating and plots is the mean of per-class MC standard deviations.
  \item Offline uncertainty plots are already wired (Step 3c).
  \item Live uncertainty-aware actuation is \emph{implemented} in \texttt{utils/inference.py} but is not currently used by the Step 7 live actuation script in this repo state.
\end{itemize}

\end{document}

